% =====================================================================
%  Distributed Multi-Head Attention with Hybrid MPI Parallelism
%  Parallel & Distributed Programming — project report (draft)
%
%  Compile:  pdflatex REPORT.tex   (run twice for the table of contents
%            and cross-references). Needs a standard TeX Live install.
%
%  Charts:   run the benchmarks on the cluster, then
%            python3 scripts/make_charts.py
%            which writes results/chart-b-n-sizing.png, chart-c-granularity.png,
%            chart-d-speedup.png. Until then, this report shows placeholder boxes
%            in their place (it still compiles).
%
%  TODO before submitting: fill every  \TODO{...}  with your measured numbers.
% =====================================================================
\documentclass[11pt,a4paper]{article}

\usepackage[margin=2.4cm]{geometry}
\usepackage{amsmath,amssymb}
\usepackage{graphicx}
\usepackage{booktabs}
\usepackage{xcolor}
\usepackage{listings}
\usepackage{float}
\usepackage{caption}
\usepackage[hidelinks]{hyperref}
\usepackage{enumitem}
\usepackage{microtype}
\sloppy  % prefer slightly loose lines over content spilling into the margin

% --- a clear, portable style for pseudocode (listings is in every TeX Live) ---
\definecolor{kw}{RGB}{0,0,160}
\definecolor{cmt}{RGB}{0,128,0}
\lstdefinestyle{pseudo}{
  basicstyle=\small\ttfamily,
  keywordstyle=\color{kw}\bfseries,
  commentstyle=\color{cmt}\itshape,
  mathescape=true,
  columns=fullflexible,
  keepspaces=true,
  frame=single,
  framesep=6pt,
  numbers=left,
  numberstyle=\tiny\color{gray},
  xleftmargin=2.4em,
  breaklines=true,
  morekeywords={for,each,do,end,if,then,else,return,in,parallel,on,all,broadcast,
    scatter,gather,allreduce,send,recv,function,Input,Output,to,from,foreach,while}
}
\lstset{style=pseudo}

\newcommand{\code}[1]{\texttt{#1}}
\newcommand{\TODO}[1]{{\color{red}\textbf{[TODO: #1]}}}

% Measured load-imbalance figures (Chart C and its granularity-adjustment table).
\newcommand{\GRANHYB}{175.7\,\%}      % hybrid G=2, N=32768, p=22 (measured)
\newcommand{\GRANTEN}{28.6\,\%}       % 1D tensor, N=32768, p=22 (measured)
% Chart D speedup at the 2X=44 oversubscription point (clean, optimized binary).
\newcommand{\Tone}{92.0}   % serial baseline T(1) at N*=32768, seconds (measured)

% Insert a chart if it exists, otherwise a placeholder box (keeps the report compilable).
\newcommand{\chartfig}[3]{% #1 = path, #2 = caption, #3 = label
  \begin{figure}[H]\centering
  \IfFileExists{#1}{\includegraphics[width=\linewidth]{#1}}{%
    \fbox{\parbox[c][4.2cm][c]{0.9\linewidth}{\centering
      \textit{Chart placeholder} --- file \code{#1} not found.\\[2pt]
      Run the benchmarks, then \code{python3 scripts/make\_charts.py}.}}}
  \caption{#2}\label{#3}\end{figure}}

\begin{document}

% =====================================================================
%  COVER PAGE
% =====================================================================
\begin{titlepage}
\centering
\vspace*{1.2cm}
{\large\scshape Hanoi University of Science and Technology}\\[0.25cm]
{\large School of Information and Communication Technology}\\[0.25cm]
\rule{0.5\linewidth}{0.4pt}\\[2.6cm]

{\large Parallel \& Distributed Programming}\\[0.35cm]
{\large\itshape Course Project Report}\\[2.4cm]

{\color{kw}\rule{\linewidth}{1.2pt}}\\[0.55cm]
{\Huge\bfseries Distributed Multi-Head Attention}\\[0.5cm]
{\Large A Hybrid (Task\,+\,Data) MPI Parallelization}\\[0.45cm]
{\color{kw}\rule{\linewidth}{1.2pt}}\\[2.4cm]

{\large\bfseries Group Members}\\[0.55cm]
{\large Nguyen Hoang Thai}\\[0.18cm]
{\large Pham Le Minh Quang}\\[0.18cm]
{\large Pham Ngoc Thach}\\[0.18cm]
{\large Dang Xuan Dat}\\

\vfill
{\large \today}
\end{titlepage}

% =====================================================================
%  ABSTRACT (own page)
% =====================================================================
\begin{abstract}
\noindent
We parallelize the Transformer \emph{multi-head attention} (MHA) operator,
$\mathrm{Attention}(Q,K,V)=\mathrm{softmax}\!\big(QK^{\top}/\sqrt{d_k}\big)V$,
implemented from scratch in C with OpenMPI (no deep-learning frameworks). Three
strategies are compared against a sequential baseline: \emph{head parallelism}
(task decomposition), \emph{tensor parallelism} (data decomposition, with a 1D
row-block path, a 2D Cannon path, and a ring/Flash-style path that streams sharded
$K,V$ through an online softmax), and a \emph{hybrid} scheme that splits the
processes into groups via \code{MPI\_Comm\_split}; within each rank an OpenMP layer
adds shared-memory parallelism. We describe the level of
parallelism, the decomposition and mapping, the communication strategy and
topology, and load balancing; give pseudocode; and report correctness plus three
experiments: execution time vs.\ input size $N$, per-process granularity / load
balance, and speedup as the process count scales. Measured on a real three-node,
$22$-core Wi-Fi cluster, all parallel paths reproduce the sequential result within
tolerance; the computation scales to about $20\times$ on $22$ cores while
communication over the wireless link caps the end-to-end speedup near $4.4\times$
--- a textbook communication wall.
\end{abstract}
\newpage

% =====================================================================
%  TABLE OF CONTENTS (own page)
% =====================================================================
\tableofcontents
\newpage

% =====================================================================
\section{Introduction}
% =====================================================================
Attention is the central operation of the Transformer architecture. For a
sequence of $N$ tokens with model dimension $d_{\text{model}}$ split across $H$
heads (so each head has width $d_k = d_{\text{model}}/H$), each head computes
\begin{equation}
\label{eq:attn}
\mathrm{Attention}(Q,K,V)=\mathrm{softmax}\!\left(\frac{QK^{\top}}{\sqrt{d_k}}\right)V,
\qquad Q,K,V\in\mathbb{R}^{N\times d_k}.
\end{equation}
The $H$ head outputs are concatenated to form the $N\times d_{\text{model}}$
result. The cost is dominated by two matrix products per head --- the
$N\times N$ score matrix $S=QK^{\top}$ and the $N\times d_k$ output $SV$ --- so a
single head costs $\mathcal{O}(N^2 d_k)$ and full MHA costs
$\mathcal{O}(H\,N^2 d_k)$. The quadratic dependence on $N$ is what makes the
problem worth distributing.

\paragraph{Background and related work.} The decompositions used here mirror the
ones that production systems use, reduced to their essentials. Splitting whole heads
across devices is \emph{head} (a form of \emph{tensor-model}) parallelism, as in
Megatron-LM; splitting a single head's matrices across devices is \emph{tensor}
parallelism. The ring path is a distributed-memory rendering of
\emph{FlashAttention}'s online-softmax recurrence (keep a running max, denominator
and weighted sum so the $N\times N$ score matrix is never materialized), combined
with the ``ring attention'' idea of streaming sharded $K,V$ around the processes;
it is also a direct cousin of the classic systolic matrix-multiply of
\emph{Cannon (1969)}, which the 2D path implements literally on an
\code{MPI\_Cart} torus. The communication primitives --- tree-based
\code{Bcast}/\code{Scatterv}/\code{Gatherv}/\code{Allreduce} and
\code{Sendrecv\_replace} ring shifts --- are the standard MPI building blocks for
these patterns. The contribution of this project is not a new algorithm but a
\emph{from-scratch, framework-free} C\,+\,MPI implementation of all of them, on one
shared kernel, measured honestly on a real heterogeneous cluster.

\paragraph{Goal.} Implement MHA in pure C\,+\,MPI, parallelize it across a small
cluster, verify correctness against a sequential reference, and characterize its
performance (runtime vs.\ size, load balance, speedup). The code is organized as
in Table~\ref{tab:files}.

\begin{table}[H]
\centering\small
\caption{Source layout.}\label{tab:files}
\begin{tabular}{@{}ll@{}}
\toprule
File & Responsibility \\
\midrule
\code{src/tensor.c}          & \code{Tensor} type: alloc/free/zero, \code{transpose} (for Cannon's $K^{\top}$) \\
\code{src/attention.c}       & streaming kernel + sequential single/multi-head attention (ground truth) \\
\code{src/data\_gen.c}       & deterministic in-memory $Q,K,V$ generation (rank 0) \\
\code{src/profiler.c}        & \code{MPI\_Wtime} timers (I/O, compute, comm, wait) + msg/byte/latency counters \\
\code{src/head\_parallel.c}  & M1 --- head parallelism (\code{MPI\_Scatterv}/\code{Gatherv}) \\
\code{src/tensor\_parallel.c}& M2 --- tensor parallelism: 1D row-block, 2D Cannon, ring (Flash-style) \\
\code{src/hybrid.c}          & M3 --- hybrid (\code{MPI\_Comm\_split} into head groups) \\
\code{src/progress.c}        & rank-0 progress/ETA meter for long runs \\
\code{src/main.c}            & CLI, dispatch, correctness check \\
\bottomrule
\end{tabular}
\end{table}

\paragraph{Roadmap.} Table~\ref{tab:roadmap} maps each question the project must
answer to the section that answers it, so the report can be read against the brief
directly.

\begin{table}[H]
\centering\small
\caption{Where each required question is addressed.}\label{tab:roadmap}
\begin{tabular}{@{}ll@{}}
\toprule
Question & Section \\
\midrule
Level of parallelism: task or data?            & \S\ref{sec:level} \\
Decomposition technique?                        & \S\ref{sec:decomp} \\
Mapping (1D / 2D $n/\sqrt{p}$ block)?            & \S\ref{sec:mapping} \\
Communication strategy \& topology; blocking?   & \S\ref{sec:comm} \\
Load balancing applied?                          & \S\ref{sec:loadbal}, \S\ref{sec:chartc} \\
Pseudocode of the parallel algorithm            & \S\ref{sec:pseudo} \\
Correctness of the parallel result              & \S\ref{sec:correctness} \\
Input size $N$ for $2$--$3$ min; time vs.\ $N$  & \S\ref{sec:chartb} (Chart~B) \\
Granularity / load balance ($25\%$ rule)        & \S\ref{sec:chartc} (Chart~C) \\
Speedup at $2N$, $p=1\ldots2X$, with/without comm & \S\ref{sec:chartd} (Chart~D), App.~\ref{app:2n} \\
\bottomrule
\end{tabular}
\end{table}

% =====================================================================
\section{Level of parallelism: task vs.\ data}
\label{sec:level}
% =====================================================================
The project deliberately implements \textbf{both} classical levels of
parallelism so they can be compared:

\begin{itemize}[leftmargin=1.4em]
\item \textbf{Task (functional) parallelism --- head mode (M1).} The $H$ heads
are mutually independent computations on disjoint slices of the columns of
$Q,K,V$. Each head is a self-contained \emph{task}; assigning whole heads to
different processes is task parallelism.
\item \textbf{Data parallelism --- tensor mode (M2).} A \emph{single} head is
distributed by partitioning its data (the rows of $Q$, or 2D blocks of the
matrices) across processes; every process runs the same code on a different
portion of the data.
\item \textbf{Hybrid (M3).} Processes are partitioned into groups; heads are
distributed across groups (task), and within each group one head at a time is
distributed by data (tensor). This is task parallelism at the outer level and
data parallelism at the inner level.
\end{itemize}

\noindent
In short: \emph{head mode = task}, \emph{tensor mode = data}, \emph{hybrid =
task over data}.

\paragraph{A third level: shared-memory threads.} Within every MPI rank the
per-query-row loop of the attention kernel is parallelized with \textbf{OpenMP}
(\code{\#pragma omp parallel for}), each thread keeping a private score buffer.
This adds a fine-grained shared-memory level beneath the MPI process level, so a
configuration is described by both the process count $p$ and the per-rank thread
count $t$ (reported as \code{threads=}$t$); with $t=1$ the runs are pure MPI.

% =====================================================================
\section{Decomposition technique}
\label{sec:decomp}
% =====================================================================
Using the standard taxonomy (data / recursive / exploratory / speculative /
hybrid), this problem is a combination of \textbf{data decomposition} and
\textbf{functional (task) decomposition}:

\begin{itemize}[leftmargin=1.4em]
\item \textbf{Data decomposition} (tensor mode). The output and intermediate
data are partitioned and each partition is computed by one process:
\begin{itemize}
  \item \emph{1D}: partition the \emph{query rows} --- process $r$ owns a
        contiguous block of rows of $Q$ and produces the corresponding rows of
        the output. Softmax is row-wise, so each row is fully local: no
        reduction is needed across processes.
  \item \emph{2D (Cannon)}: partition all matrices into a
        $\sqrt{p}\times\sqrt{p}$ grid of blocks; one block per process.
\end{itemize}
\item \textbf{Functional/task decomposition} (head mode). The computation is
split by \emph{function instances} --- one independent attention computation per
head.
\item \textbf{Hybrid} (hybrid mode). Both are combined.
\end{itemize}

\noindent
The remaining categories do not apply: there is no recursive subdivision of the
problem (\emph{recursive}), no search space explored in parallel
(\emph{exploratory}), and no redundant ``guess-ahead'' work that may be discarded
(\emph{speculative}).

% =====================================================================
\section{How the parallelization is implemented}
% =====================================================================

\subsection{Implementation overview and data flow}
\label{sec:impl}
All three modes share one control flow, in \code{main.c}: parse and \emph{validate}
the configuration collectively (every rank parses the same \code{argv}, so all reach
the same verdict and abort together on a bad configuration --- no rank ever proceeds
into a kernel with, say, a non-divisible process count); generate the input;
dispatch to the chosen kernel; then, on rank~0, recompute the result sequentially
and report the maximum absolute error. The dispatch is a thin switch over
\code{MODE\_\{HEAD,TENSOR,HYBRID\}}; the heavy lifting lives in one file per mode
(Table~\ref{tab:files}), which keeps the communication strategy of each mode
isolated and independently testable.

\paragraph{Deterministic input without a parallel file system.}
The cluster has no shared file system, so rather than read $Q,K,V$ from disk we
\emph{generate} them: \code{data\_gen\_fill} seeds \code{srand} with a fixed
per-matrix constant and fills the tensor with values in $[-1,1]$. Seeding makes the
input bit-identical on every run, so the correctness check and the timings are
reproducible. Only rank~0 needs the full matrices (it is the master that scatters
them), so after the memory optimization (Section~\ref{sec:discussion}) the other
ranks allocate them empty and receive their slices over MPI.

\paragraph{One streaming kernel, reused everywhere.}
The numerically stable, single-query-row routine \code{attention\_row} ---
$o=\mathrm{softmax}(qK^{\top}/\sqrt{d_k})\,V$ for one query $q$ --- is the
\emph{only} compute kernel in the project. It processes one score row at a time
with $\mathcal{O}(N)$ scratch instead of forming the $N\times N$ score matrix, so
its memory is linear in $N$ and its arithmetic order matches the reference exactly
(hence the bit-exact 1D result). The sequential baseline, the 1D tensor path and
the inner loop of hybrid all call it on their local query rows; the ring path
inlines an \emph{online} variant of the same recurrence so it can fold in one
$K,V$ block at a time. Because every path funnels through this kernel, correctness
established once carries across all modes, and an OpenMP \code{parallel for} over
the query-row loop adds a shared-memory level for free (each thread keeps a private
score buffer).

\paragraph{Self-describing timing.}
Each kernel brackets its phases with the profiler
(Section~\ref{sec:profiling}), and rank~0 prints either a human summary or a
per-rank CSV line. No external profiler or trace tool is needed; the program
explains its own time and communication, which is what makes every chart in this
report reproducible from a single binary.

\subsection{Mapping: process / processor assignment}
\label{sec:mapping}

\paragraph{Head mode (1D over heads).}
With $H$ heads and $p$ processes, process $r$ receives
$\lfloor H/p\rfloor$ consecutive heads, and the last rank takes the remainder
$H \bmod p$. The root scatters the corresponding $d_k$-wide column slices of each
row of $Q,K,V$ with \code{MPI\_Scatterv} and collects the results with
\code{MPI\_Gatherv}.

\paragraph{Tensor mode, 1D (row-block).}
The $N$ query rows are mapped to $p$ processes in contiguous blocks. With
$\text{base}=\lfloor N/p\rfloor$ and $\text{rem}=N\bmod p$, the first
$\text{rem}$ ranks own $\text{base}+1$ rows and the rest own $\text{base}$ rows,
so any two processes differ by at most one row. Each process holds its block of
$Q$ (scattered) and \emph{full} $K,V$ (broadcast), and computes its output rows
entirely locally.

\paragraph{Tensor mode, 2D (Cannon).}
The processes form a $\sqrt{p}\times\sqrt{p}$ Cartesian grid
(\code{MPI\_Cart\_create} with periodic dimensions). Each process owns an
$\frac{N}{\sqrt{p}}\times\frac{d_k}{\sqrt{p}}$ block of $Q$ (and the matching
blocks of $K^{\top}$ and $V$) --- exactly the
$n/\sqrt{p}\times n/\sqrt{p}$ block mapping. This path requires $p$ to be a
perfect square and $N,d_k$ divisible by $\sqrt{p}$.

\paragraph{Tensor mode, ring (Flash-style).}
The $N$ query rows \emph{and} the $N$ key/value rows are both partitioned into
contiguous blocks of $N/p$, so rank $r$ owns its query block and only its own
$N/p$ rows of $K,V$ (not the whole matrices). The $K,V$ blocks are streamed around
a ring of the $p$ ranks with non-blocking \code{MPI\_Isend}/\code{MPI\_Irecv}
double buffering; at each of the $p$ steps a rank folds the block it currently
holds into a running, numerically stable \emph{online softmax} over its local
queries (the Flash-attention update: keep a running max $m$, running denominator
$\ell$, and running weighted sum, rescaling on each new max). It therefore never
materializes a score row and never holds more than $O(N/p)$ of $K,V$, and the
shift overlaps the compute. Works for any $p$; the result matches the baseline up
to floating-point reassociation ($\sim10^{-7}$).

\paragraph{Hybrid mode (task $\times$ data).}
\code{MPI\_Comm\_split} partitions \code{MPI\_COMM\_WORLD} into
$G$ groups of $p/G$ processes (the colour is $\lfloor r/(p/G)\rfloor$). The $H$
heads are split so each group owns $H/G$ heads (task level); within a group the
$p/G$ processes run the 1D tensor-parallel kernel on one head at a time (data
level). The validation in \code{main.c} enforces $p \bmod G = 0$ and
$H \bmod G = 0$ so the process grid is rectangular --- e.g.\ with $p=22$, $G=2$,
$H=4$ each group is $11$ processes computing $2$ heads.

\subsection{Communication strategy and topology}
\label{sec:comm}

\paragraph{Master--slave (star).}
Rank~0 is the master. It owns the full $Q,K,V$, distributes work, and collects
results. In head and tensor mode this uses the collective
\code{Scatterv}/\code{Bcast}/\code{Gatherv}; in hybrid mode the master distributes
each group's slice to the group roots with \code{MPI\_Scatter} and collects the
results with \code{MPI\_Gather} over a dedicated group-roots sub-communicator
(replacing an earlier serial \code{Send}/\code{Recv} fan-out). Logically this is a
star topology centred on rank~0.

\paragraph{Blocking vs.\ non-blocking.}
The head, 1D-tensor, Cannon and hybrid paths use \textbf{blocking} collectives and
point-to-point calls, which keeps them simple and correct at the cost of no
compute/communication overlap. The \textbf{ring} path instead uses
\textbf{non-blocking} \code{MPI\_Isend}/\code{MPI\_Irecv} closed by an
\code{MPI\_Waitall}, deliberately overlapping each $K,V$ shift with the
online-softmax compute of the block currently held (Section~\ref{sec:discussion}).

\paragraph{Topologies in use.}
\begin{itemize}[leftmargin=1.4em]
\item \textbf{Tree} --- the MPI collectives (\code{Bcast}, \code{Scatterv},
\code{Gatherv}, \code{Allreduce}) are realized by the library with tree-based
algorithms, giving $\mathcal{O}(\log p)$ depth.
\item \textbf{Ring / 2D torus} --- the Cannon path builds a periodic 2D
Cartesian communicator and performs the two matrix multiplies as a sequence of
$\sqrt{p}$ ring shifts using \code{MPI\_Sendrecv\_replace} (rows shift left,
columns shift up). This is the classic torus communication pattern.
\item \textbf{Sub-communicator reduction} --- in the Cannon path a single row of
the score matrix is split across a block-row, so the numerically stable softmax
is completed with two \code{MPI\_Allreduce} calls (max, then sum) over a
\emph{row} sub-communicator created by \code{MPI\_Comm\_split}.
\end{itemize}

\noindent
Table~\ref{tab:comm} summarizes the primitives per mode.
\begin{table}[H]
\centering\small
\caption{Communication primitives by mode.}\label{tab:comm}
\begin{tabular}{@{}llll@{}}
\toprule
Mode & Distribute & Reduce/exchange & Collect \\
\midrule
Head (M1)        & \code{Scatterv}                 & ---                                & \code{Gatherv} \\
Tensor 1D (M2)   & \code{Scatterv}(Q), \code{Bcast}(K,V) & ---                          & \code{Gatherv} \\
Tensor Cannon    & 2D block scatter (\code{Send})  & \code{Sendrecv\_replace}, \code{Allreduce} & 2D block gather \\
Tensor Ring      & \code{Scatterv}(Q,K,V)          & \code{Isend}/\code{Irecv} ring shifts & \code{Gatherv} \\
Hybrid (M3)      & \code{Comm\_split}, \code{Scatter} + inner M2 & inner M2          & \code{Gather} \\
\bottomrule
\end{tabular}
\end{table}

\subsection{Load balancing}
\label{sec:loadbal}
Two kinds of balance must be distinguished: \emph{static} balance (does the
mapping hand every process an equal \emph{amount} of work?) and \emph{realized}
balance (do the processes actually \emph{finish} together?). The mappings here are
statically balanced:
\begin{itemize}[leftmargin=1.4em]
\item \textbf{Tensor 1D} is balanced by construction: the remainder rows are
spread one-each over the first $\text{rem}$ ranks, so process workloads differ by
at most one row of work --- negligible for large $N$.
\item \textbf{Head mode} assigns $\lfloor H/p\rfloor$ heads per rank and gives
the remainder to the \emph{last} rank, so when $H\bmod p\neq 0$ the last rank does
strictly more work; this is the structural source of imbalance in M1.
\item \textbf{Hybrid} gives every group equal processes and equal heads when
$p\bmod G=0$ and $H\bmod G=0$ (enforced in \code{main.c}).
\end{itemize}
Static balance, however, assumes \emph{homogeneous} processes and \emph{free}
communication --- neither of which holds on our cluster. The three CPU generations
run the same row block at different speeds, the Wi-Fi link is slow and variable,
and the group roots carry the distribution traffic. Realized balance is therefore
an experimental question: the granularity benchmark (Section~\ref{sec:chartc})
measures each process's compute$+$comm time and flags any pair differing by more
than $25\%$. As Chart~C shows, the statically ``balanced'' hybrid mapping is in
fact $175.7\%$ imbalanced in practice, and the fix is to change the decomposition
(coarser/finer), not just the chunk count.

\subsection{Cost analysis (computation vs.\ communication)}
\label{sec:cost}
Table~\ref{tab:cost} gives the per-process asymptotic costs, which explain the
experimental behaviour. For one head of width $d_k$ on $p$ processes, the
sequential work is $\Theta(N^2 d_k)$; all data paths cut the \emph{computation}
to $\Theta(N^2 d_k/p)$ but differ sharply in \emph{communication volume} and in
the \emph{memory} each rank must hold --- the quantity that decides whether the
algorithm is compute- or communication-bound on a given interconnect.

\begin{table}[H]
\centering\small
\caption{Per-process cost for one head ($N$ tokens, width $d_k$, $p$ processes).
Communication volume is the bytes a rank sends/receives; ``$K,V$ resident'' is the
peak key/value memory it must store.}\label{tab:cost}
\begin{tabular}{@{}lccc@{}}
\toprule
Path & Computation & Communication volume & $K,V$ resident \\
\midrule
Sequential        & $\Theta(N^2 d_k)$      & ---                              & $\Theta(N d_k)$   \\
Head (task)       & $\Theta(N^2 d_k\,H/p)$ & $\Theta(N d_k\,H/p)$            & $\Theta(N d_k)$   \\
Tensor 1D         & $\Theta(N^2 d_k/p)$    & $\Theta(N d_k)$ (full $K,V$ bcast) & $\Theta(N d_k)$   \\
Tensor Cannon     & $\Theta(N^2 d_k/p)$    & $\Theta(N d_k/\sqrt{p})$ per shift & $\Theta(N d_k/p)$ \\
Tensor Ring       & $\Theta(N^2 d_k/p)$    & $\Theta(N d_k)$ ($p{-}1$ shifts)   & $\Theta(N d_k/p)$ \\
\bottomrule
\end{tabular}
\end{table}

\noindent
Two consequences drive the results. First, in 1D and ring the computation falls
as $1/p$ but the communication volume \emph{does not} (every rank still needs all
of $K,V$ once), so the communication-to-computation ratio grows like $p$: past
some $p$ the broadcast/streaming dominates --- exactly the ``communication wall''
the speedup curve hits on Wi-Fi (Section~\ref{sec:chartd}). Second, only Cannon
and ring shrink the resident $K,V$ to $\Theta(N d_k/p)$; 1D replicates the full
$K,V$ on every rank, which is what ultimately caps $N$ on a memory-limited node
and motivates the sharded ring path (Chart~E).

\paragraph{A worked estimate of the crossover.} Plugging in our numbers makes the
regime concrete. For one head, $d_k=128$, the broadcast moves $K$ and $V$, i.e.\
$2Nd_k$ floats $=2N\!\cdot\!128\!\cdot\!4$ bytes; the per-rank compute is
$\Theta(N^2 d_k/p)$ flops. At the small end ($N=2048$, $p=22$) the compute is only
$\sim\!2{\cdot}2048^2{\cdot}128/22\approx4.9\times10^{7}$ flops --- microseconds on a
CPU --- against a $\sim\!2$\,MB broadcast that, at $\sim\!1$\,MB/s shared Wi-Fi,
takes seconds: communication outweighs computation by orders of magnitude, which is
exactly the $97$--$98\%$ communication share Chart~B reports below $N=8192$. The
compute term grows quadratically while the broadcast grows only linearly, so the two
cross over near $N\approx32768$ (the measured share there falls to $\approx\!55\%$),
and only at $N\gtrsim131072$ would compute finally dominate --- but such a run takes
$\sim\!30$ minutes, outside the $2$--$3$ minute budget. This single back-of-envelope
calculation predicts every qualitative feature of Charts~B and~D.

\subsection{Pseudocode}
\label{sec:pseudo}

\noindent\textbf{Sequential reference (per head).}
\begin{lstlisting}
function ATTENTION(Q, K, V):            # Q,K,V : N x d_k
    S $\gets$ Q $\cdot$ K$^\top$                      # N x N scores
    S $\gets$ S / $\sqrt{d_k}$
    for each row i of S:                # numerically stable softmax
        m $\gets$ max_j S[i,j]
        S[i,:] $\gets$ exp(S[i,:] - m); S[i,:] $\gets$ S[i,:] / sum(S[i,:])
    return S $\cdot$ V                       # N x d_k
\end{lstlisting}

\noindent\textbf{M2 --- 1D tensor-parallel attention} (the core parallel kernel;
also used inside hybrid).
\begin{lstlisting}
function TP_ATTENTION_1D(Q, K, V) on communicator of p ranks:
    broadcast N, d_k from rank 0
    base $\gets$ N / p;  rem $\gets$ N mod p
    my_rows $\gets$ base + (rank < rem ? 1 : 0)        # 1D row mapping
    # --- communication ---
    scatter row-blocks of Q  -> Q_local               # MPI_Scatterv
    broadcast full K, V to every rank                 # MPI_Bcast
    # --- computation (entirely local; softmax is per-row) ---
    for each local row i:
        for j in 0..N-1: s[j] $\gets$ (Q_local[i] $\cdot$ K[j]) / $\sqrt{d_k}$
        softmax(s)                                    # local: row is not split
        O_local[i] $\gets$ s $\cdot$ V                          # row-major accumulation
    # --- communication ---
    gather O_local -> O on rank 0                     # MPI_Gatherv
    return O
\end{lstlisting}

\noindent\textbf{M2 --- 2D Cannon path} (perfect-square $p$).
\begin{lstlisting}
function TP_ATTENTION_CANNON(Q, K, V) on $\sqrt{p}\times\sqrt{p}$ torus:
    scatter Q, K$^\top$, V into $\sqrt{p}\times\sqrt{p}$ blocks
    S $\gets$ CANNON_MATMUL(Q_block, Kt_block)         # ring shifts, MPI_Sendrecv_replace
    S $\gets$ S / $\sqrt{d_k}$
    DISTRIBUTED_SOFTMAX(S over row sub-comm)          # MPI_Allreduce max, then sum
    O $\gets$ CANNON_MATMUL(S_block, V_block)
    gather O blocks -> rank 0
\end{lstlisting}

\noindent\textbf{M2 --- ring (Flash-style) path} (sharded $K,V$, overlapped).
\begin{lstlisting}
function TP_ATTENTION_RING(Q, K, V) on a ring of p ranks:
    scatter row-blocks of Q, K, V -> Q_loc, K_blk, V_blk    # each rank: N/p rows
    m[i] $\gets -\infty$;  l[i] $\gets$ 0;  acc[i] $\gets$ 0          # online-softmax state
    for step in 0..p-1:
        if step < p-1: Isend(K_blk,V_blk -> next); Irecv(<- prev)   # overlap shift
        foreach local query i, each key j in current block:    # parallel-for (OpenMP)
            s $\gets$ (Q_loc[i] $\cdot$ K_blk[j]) / $\sqrt{d_k}$
            if s > m[i]: rescale acc[i],l[i] by exp(m[i]-s); m[i] $\gets$ s   # new max
            p $\gets$ exp(s - m[i]); l[i] += p; acc[i] += p * V_blk[j]
        if step < p-1: Waitall; swap current <- received block
    O_loc[i] $\gets$ acc[i] / l[i]
    gather O_loc -> rank 0
\end{lstlisting}

\noindent\textbf{M3 --- Hybrid.}
\begin{lstlisting}
function HYBRID_MHA(Q, K, V, H heads, G groups) on p ranks:
    split COMM_WORLD into G groups of p/G ranks         # MPI_Comm_split
    build a sub-communicator of the G group roots       # MPI_Comm_split
    master scatters each group its H/G head-columns     # MPI_Scatter (roots comm)
    for each head h owned by my group:                  # task level
        O_h $\gets$ TP_ATTENTION_1D(Q_h, K_h, V_h) on my group   # data level
    master gathers results from the group roots         # MPI_Gather (roots comm)
\end{lstlisting}

% =====================================================================
\section{Profiling methodology}
\label{sec:profiling}
% =====================================================================
Timing uses \code{MPI\_Wtime} through a small profiler that accumulates four time
categories per rank --- \textbf{I/O} (input generation), \textbf{compute} (local
matrix work), \textbf{communication} (time inside MPI data-transfer calls), and
\textbf{wait} (idle time blocked at a synchronization point) --- together with two
communication-volume counters, \textbf{msgs} (number of MPI data transfers the rank
took part in) and \textbf{bytes} (volume it sent/received), and a one-off ping-pong
probe that measures the cluster's worst-case one-way \textbf{link latency}. Each
parallel kernel brackets its phases with
\code{profiler\_start}/\code{profiler\_stop}, so the ``with communication'' time is
$t_{\text{compute}}+t_{\text{comm}}$ and the ``without communication'' time is
$t_{\text{compute}}$ alone. Wait is separated from real transfer by bracketing an
\code{MPI\_Barrier} into the wait timer just before each collective (enabled with
\code{--profile-wait}; off by default so headline timing runs are unperturbed). With
\code{--csv} the program prints one line per rank
(\code{rank,seq\_len,t\_io,t\_compute,t\_comm,t\_wait,msgs,bytes,latency\_us}); the
benchmark scripts collect these into \code{results/*.csv}, which
\code{scripts/make\_charts.py} turns into the figures below --- Chart~C stacks
compute/comm/wait per rank and annotates message and byte totals, and Chart~D adds a
communication-volume and link-latency panel versus $p$. Long runs can be watched
live with \code{--progress}, which prints a rank-0 percentage and ETA to
\code{stderr}.

% =====================================================================
\section{Results}
% =====================================================================

\subsection{Experimental platform}
\label{sec:platform}
All measurements were taken on a real \textbf{three-node cluster} of commodity
laptops connected over a shared \textbf{Wi-Fi LAN} (subnet
\code{192.168.124.0/24}); the nodes are named \code{master}, \code{slave1} and
\code{slave2} and communicate by host name through \code{/etc/hosts}. Every node
runs the \emph{same} software stack --- Ubuntu~24.04.4, glibc~2.39 and
\textbf{Open~MPI~4.1.6} --- so a single binary built on \code{master} runs
unchanged on all three; passwordless SSH lets \code{mpirun} launch ranks on the
slaves, and the binary is staged to an identical path
(\code{/tmp/mpiattn/hybrid\_attention}) on each node. The hardware is summarized
in Table~\ref{tab:cluster}. It is deliberately \emph{heterogeneous} --- three
different CPU generations and per-core speeds --- which matters for the
load-balance analysis in Section~\ref{sec:chartc}.

\begin{table}[H]
\centering\small
\caption{Cluster hardware. ``Phys.'' = physical cores used as MPI slots;
``Log.'' = logical (SMT) threads available to OpenMP.}\label{tab:cluster}
\begin{tabular}{@{}llccc@{}}
\toprule
Node & CPU & Phys. & Log. & Max clock \\
\midrule
\code{master} & AMD Ryzen 7 7840H        & 8  & 16 & 5.14\,GHz \\
\code{slave1} & Intel Core i5-1135G7     & 4  & 8  & 4.20\,GHz \\
\code{slave2} & Intel Core i7-13620H     & 10 & 16 & 4.90\,GHz \\
\midrule
\textbf{Total} & --- & \textbf{22} & \textbf{40} & --- \\
\bottomrule
\end{tabular}
\end{table}

\noindent
The total physical-core count is therefore $X=22$; following the brief, the
process count is set equal to the number of physical cores ($p=22$) for the
size and granularity experiments, and swept up to $2X=44$ for the speedup
experiment. Unless noted, \code{OMP\_NUM\_THREADS=1} so the figures show
\emph{pure MPI} scaling; the inter-node link is measured by a ping-pong probe at
a one-way latency of roughly \textbf{4--6\,ms} (typical of contended Wi-Fi),
which --- as the results show --- makes the cluster strongly
\emph{communication-bound}.

\paragraph{A note on the network regime.} Because the interconnect is Wi-Fi
rather than wired Ethernet or InfiniBand, communication is two to three orders of
magnitude slower than the CPUs. This is the dominant feature of every result
below: it sharpens the contrast between the ``with-communication'' and
``computation-only'' curves the brief asks for, and it makes the
shared-memory-to-network transition directly visible in the speedup curve
(Section~\ref{sec:chartd}). We report it honestly rather than hiding it, because
characterizing communication cost is one of the graded objectives.

\paragraph{Measured network characteristics.} The built-in ping-pong probe
(Section~\ref{sec:profiling}) reports a one-way latency that rises from about
$0.3\,\mu$s for two ranks sharing \code{master} (pure shared memory) to
$\mathbf{4}$--$\mathbf{7.5}$\,ms once ranks span the wireless LAN --- a four-order-of-magnitude
gap that already foretells the communication wall. Effective bandwidth is
correspondingly low and, crucially, \emph{shared}: Wi-Fi is a half-duplex broadcast
medium, so the per-rank throughput of a collective falls as more ranks contend for
the air. From the 1D run at $p=16$, $N=16384$ (Chart~E) the full $K,V$ broadcast of
$\approx\!17$\,MB took $13.8$\,s, an effective $\sim\!1.2$\,MB/s per rank under
$16$-way contention --- roughly three orders of magnitude below the $\sim\!10$\,GB/s
the CPUs sustain on local memory. It is this ratio, not any algorithmic flaw, that
makes an $\mathcal{O}(N^2)$-compute / $\mathcal{O}(N)$-communication operator behave
as communication-bound at the sizes the $2$--$3$ minute budget permits.

\subsection{Correctness verification}
\label{sec:correctness}
\textbf{Methodology.} ``Does the parallel program compute the right answer?'' is
settled by an \emph{oracle} comparison built into every run. The same deterministic
$Q,K,V$ that the parallel kernel consumes are also fed to the independent sequential
routine \code{mha\_seq} on rank~0; this sequential routine is the textbook
definition of attention (Eq.~\ref{eq:attn}) with no parallelism, so it is a trusted
ground truth. We then report the element-wise maximum absolute difference
$\max_i|o^{\text{par}}_i-o^{\text{seq}}_i|$ and pass the run if it is below
$10^{-3}$ (single-precision attention outputs are $O(1)$, so $10^{-3}$ is a strict
relative tolerance of $\sim\!0.1\%$, yet loose enough to admit benign floating-point
reordering). Because all parallel paths funnel through the \emph{same}
single-row kernel as the baseline (Section~\ref{sec:impl}), a pass certifies the
distribution/collection logic specifically, with the arithmetic held constant.

\textbf{Results.} The 1D, head and hybrid paths reproduce the baseline
\emph{bitwise} ($\max\text{err}=0$): they not only compute the same values but
accumulate them in the same order. The Cannon and ring paths differ by only
$\sim\!10^{-7}$ --- they are mathematically identical but sum the keys in a
different order (Cannon by 2D blocks, ring by streamed shards), and finite-precision
addition is not associative, so this round-off is expected, not a bug; it is six
orders of magnitude inside tolerance. Table~\ref{tab:correct} lists the verdicts
measured on the cluster at $N=1024$, $d_{\text{model}}=512$, $H=4$, and the same
checks pass at every size used in the experiments (the benchmark scripts pass
\code{--no-check} only to keep the \emph{timed} runs clean, never to skip validation
of the configuration).

\begin{table}[H]
\centering\small
\caption{Correctness on the cluster. All modes pass; the 1D/head/hybrid paths are
bit-exact, the reordered (Cannon, ring) paths differ only at single-precision
round-off.}\label{tab:correct}
\begin{tabular}{@{}lllcl@{}}
\toprule
Mode & Path & Procs & \code{max\_err} & Verdict \\
\midrule
\code{head}   & task, \code{Scatterv}/\code{Gatherv}       & 4  & \code{0.00e+00} & PASS (exact) \\
\code{tensor} & data, 1D row-block                          & 22 & \code{0.00e+00} & PASS (exact) \\
\code{tensor} & data, ring (Flash, online softmax)          & 22 & \code{1.86e-07} & PASS \\
\code{tensor} & data, Cannon (2D torus)                     & 16 & \code{1.34e-07} & PASS \\
\code{hybrid} & task\,$\times$\,data, \code{Comm\_split}    & 22 & \code{0.00e+00} & PASS (exact) \\
\bottomrule
\end{tabular}
\end{table}

\noindent
A representative line of program output is:
\begin{center}\small\ttfamily
[HYBRID] wall=3.0801s\ \ seq\_len=1024\ \ d\_model=512\ \ heads=4\ \ procs=22\\
CORRECTNESS PASS\quad max\_err=0.00e+00
\end{center}
\noindent
The parallel program therefore accurately represents the solution to the problem
across all three modes and all four tensor-path variants.

\subsection{Chart B --- execution time vs.\ input size $N$}
\label{sec:chartb}
\textbf{Setup.} Processes $p=X=22$ (all physical cores), hybrid mode. Sweep
$N\in\{64,128,\dots,32768\}$ and plot the execution time of the slowest rank
(the ``start to finish'' time), once \emph{with} communication
($\max_r t^r_{\text{compute}}+t^r_{\text{comm}}$) and once \emph{without}
($\max_r t^r_{\text{compute}}$). The chosen size $N^{\star}$ is the one whose
runtime lands in the $2$--$3$ minute window; it is reused in the next two
experiments. Generated with \code{scripts/bench\_n\_sizing.sh}.

\chartfig{results/chart-b-n-sizing.png}
  {Execution time vs.\ input size $N$ at $p=22$, with communication (blue) and
   computation-only (orange). Log-scaled $N$ axis; the shaded band is the
   $2$--$3$ minute target.}{fig:b}

\noindent\textbf{Observation --- we choose $N^{\star}=32768$.} Table~\ref{tab:b}
lists the two curves. The computation-only time scales cleanly as
$\mathcal{O}(N^2)$ --- each doubling of $N$ multiplies it by $\approx\!4$
($2.07\to8.83\to37.8$\,s for $N=8192\to16384\to32768$) --- confirming the kernel
is correct and that the data decomposition divides the quadratic work as
expected. The with-communication curve sits far above it: communication is
\textbf{the dominant cost} over Wi-Fi, taking $87$--$98\,\%$ of the runtime up to
$N=16384$. Only at $N=32768$ does compute grow enough to pull the communication
share down to $\approx\!55\,\%$, i.e.\ the algorithm finally approaches the
compute-bound regime. At $N^{\star}=32768$ the end-to-end run takes about
\textbf{1.5--2.5 minutes} at $p=22$ (the synchronization-instrumented Chart~C run
of the same size measures $\approx\!147$\,s wall, squarely in the target band),
so $N^{\star}=32768$ is used for the granularity and (doubled to $65536$)
speedup experiments.

\begin{table}[H]
\centering\small
\caption{Chart B data (slowest rank). The gap between the two columns is the
communication overhead; its share of the total falls as $N$ grows because compute
is $\mathcal{O}(N^2)$ while the $K,V$ traffic is only $\mathcal{O}(N)$.}\label{tab:b}
\begin{tabular}{@{}rrrr@{}}
\toprule
$N$ & Computation only (s) & With communication (s) & Comm.\ share \\
\midrule
2048   & 0.17  & 6.47  & 97\,\% \\
4096   & 0.53  & 16.14 & 97\,\% \\
8192   & 2.07  & 22.40 & 91\,\% \\
16384  & 8.83  & 68.88 & 87\,\% \\
\textbf{32768} & \textbf{37.76} & \textbf{83.35} & \textbf{55\,\%} \\
\bottomrule
\end{tabular}
\end{table}

\subsection{Chart C --- granularity and load balance}
\label{sec:chartc}
\textbf{Setup.} $N=N^{\star}=32768$, $p=X=22$ (one rank per physical core),
hybrid mode with $G=2$ groups, \code{--profile-wait} on so idle time is split out
as a third (red) segment. Each bar is one process; the stack shows its
\emph{computation}, \emph{communication} and \emph{wait} time. The system is
declared balanced only if no two processes' totals differ by more than $25\%$.
Generated with \code{scripts/bench\_granularity.sh}; the script and
\code{make\_charts.py} both print the realized imbalance $(\max-\min)/\min$.

\chartfig{results/chart-c-granularity.png}
  {Per-process computation / communication / wait time at $N=N^{\star}=32768$,
   $p=22$, hybrid $G=2$. Ranks 0--7 run on \code{master}, 8--11 on \code{slave1},
   12--21 on \code{slave2}.}{fig:c}

\noindent\textbf{Observation --- the system is \emph{not} balanced.} The measured
imbalance is \textbf{175.7\,\%}, far above the $25\,\%$ threshold. The fastest
process (rank~1 on \code{master}) finishes its compute+comm in $53.4$\,s while the
slowest (rank~11, the group-1 root on \code{slave1}) takes $147.2$\,s. Three
causes combine, all visible in the chart:
\begin{enumerate}[leftmargin=1.6em]
\item \textbf{Group roots carry the distribution traffic.} The two group roots
(ranks~0 and~11) each move $207$\,MB in $10$ messages, versus $73$\,MB in $8$
messages for an ordinary member --- a $3\times$ communication load --- because all
of the master$\leftrightarrow$group transfer funnels through them.
\item \textbf{Heterogeneous compute.} Identical row counts cost $\approx\!36$\,s on
a \code{master} rank but only $6$--$15$\,s on a \code{slave2} rank: the three CPU
generations and their memory systems run the bandwidth-bound kernel at very
different rates, so an even row split is \emph{not} an even time split.
\item \textbf{Wi-Fi communication variance.} Communication ranges from $18$\,s for
a \code{master}-local rank to $133$\,s for the \code{slave1} group root over the
shared wireless link.
\end{enumerate}
The wait segments make the consequence explicit: the fast \code{master} ranks sit
idle for up to $83$\,s at the synchronization barriers, waiting for the
communication-bound straggler.

\paragraph{Adjusting the granularity.} Following the brief, because the imbalance
exceeds $25\,\%$ we adjust it. The decisive factor here is \emph{communication},
not compute-chunk size, so the effective lever is the decomposition that removes
the group-root bottleneck: switching from hybrid ($G=2$, two serializing roots) to
the pure 1D data-parallel kernel, in which $K,V$ are distributed by a balanced
\code{MPI\_Bcast} \emph{tree} and every rank owns an equal $N/p$ row block with no
privileged root. Table~\ref{tab:gran-adjust} shows the effect is dramatic: the same
problem in 1D tensor mode drops the imbalance from $175.7\%$ to \textbf{28.6\,\%}
--- a $6\times$ reduction that lands just above the $25\%$ target. The small residual
is the irreducible hardware heterogeneity (the slowest \code{slave1} core finishes
its equal row block later than the fastest \code{master} core: $60.5$\,s vs
$77.9$\,s) plus Wi-Fi jitter; closing the last few percent would require
\emph{weighted} distribution (give faster nodes proportionally more rows) or the
sharded ring path (Section~\ref{sec:discussion}), both beyond the static,
equal-split granularity the brief considers. The practical conclusion is that the
\emph{decomposition choice}, not the chunk size, is the effective granularity lever
on this cluster.

\begin{table}[H]
\centering\small
\caption{Granularity adjustment at $N=32768$, $p=22$: moving from the
two-group hybrid decomposition to the rootless 1D row-block decomposition lowers
the load imbalance below the structural-bottleneck level.}\label{tab:gran-adjust}
\begin{tabular}{@{}lcl@{}}
\toprule
Decomposition & Imbalance $(\max-\min)/\min$ & Verdict \\
\midrule
Hybrid, $G=2$ (two group roots) & \GRANHYB & \\
1D tensor row-block (no roots)  & \GRANTEN & \\
\bottomrule
\end{tabular}
\end{table}

\subsection{Chart D --- speedup}
\label{sec:chartd}
\textbf{Setup.} Vary $p=1,2,4,8,16,22,32,44$ --- the powers of two plus $X=22$
(physical cores) and $2X=44$ (full oversubscription). We measure speedup in
\textbf{tensor mode}, because a speedup curve requires the \emph{same} kernel at
every $p$ including the $p=1$ anchor, and the hybrid mode cannot run at $p=1$ (it
needs $p\ge G$). $T(1)$ is measured with the same profiler as $T(p)$, so the
with/without-communication speedups are apples-to-apples; the entire sweep uses one
(memory-optimized) binary so the points are mutually consistent. We run this study
at $N=N^{\star}=32768$ rather than $2N^{\star}=65536$: at $65536$ every point
exceeds ten minutes of wall time and, more importantly, the single-sampled Wi-Fi
communication time varies by more than $2\times$ between repeats (e.g.\ the $p=22$
broadcast was measured at both $97$\,s and $226$\,s), which would dominate the curve
with network noise rather than scaling behaviour; the $\mathcal{O}(N^2)/\mathcal{O}(N)$
compute/communication structure --- and hence the scaling --- is identical at
$N^{\star}$. Generated with \code{scripts/bench\_speedup.sh}.

\chartfig{results/chart-d-speedup.png}
  {Left: runtime vs.\ $p$ (with / without comm). Right: speedup $T(1)/T(p)$ vs.\
   $p$, against the ideal linear line.}{fig:d}

\noindent\textbf{Observation --- computation scales; communication caps it.}
Table~\ref{tab:d} reports both speedups (rank-0 timings, same profiler at every
$p$). The two curves tell opposite stories and the split between them is the key
result. \emph{Up to $p=8$ all ranks land on \code{master}} (its slot count is
$8$), so there is no network: the with- and without-communication curves coincide
and climb together to about $4.4\times$, flattening at $p=4$--$8$ as the ranks
begin to share \code{master}'s memory bandwidth (the same wall quantified in
Chart~F). \emph{From $p=16$ the job spans all three nodes}, and the curves split
sharply. The \emph{computation-only} speedup leaps --- the slaves add independent
memory systems --- to $14.9\times$ at $p=16$ and about $\mathbf{20\times}$ at the
$22$ physical cores, climbing further only as an artifact of root sampling (see
below). But the \emph{with-communication} speedup \emph{collapses} the moment the
network appears: it falls from $4.4\times$ at $p=8$ to $2.45\times$ at $p=16$ and
$1.62\times$ at $p=22$, because the $K,V$ broadcast and the result gather now cross
the Wi-Fi link, which costs more than the extra cores save. The single most telling
number in the whole study is therefore this: \textbf{the best end-to-end speedup,
$4.4\times$, is reached at $p=8$ on one machine; adding the other two machines makes
the computation up to $20\times$ faster yet leaves the wall-clock \emph{slower},
because communication dominates.} This is the communication wall in its starkest
form, and it is exactly the with- vs.\ without-communication distinction the brief
asks for. Two caveats keep the reading honest. (i) Because the table samples rank~0,
which holds fewer query rows as $p$ grows, the computation-only \emph{number} keeps
rising past $p=22$ ($25\times$ at $p=44$); this reflects root-process sampling, not
super-linear oversubscription --- the real message is simply that compute has ceased
to be the bottleneck. (ii) The partial recovery of the with-communication speedup at
$p=32$--$44$ ($3.0$--$3.5\times$) comes from oversubscription packing more ranks
back onto the high-core nodes, shortening the broadcast's Wi-Fi path; it still never
regains the single-node $p=8$ peak. In Amdahl terms, $80$--$90\%$ of the cross-node
runtime is unparallelizable transfer (e.g.\ $52$\,s of the $57$\,s at $p=22$),
which bounds the end-to-end speedup to low single digits no matter how many cores
are added.

\begin{table}[H]
\centering\small
\caption{Chart D speedup at $N^{\star}=32768$, $T(1)=\Tone$\,s. ``No comm'' uses
$t_{\text{compute}}$ only; ``with comm'' uses $t_{\text{compute}}+t_{\text{comm}}$.
$p\le8$ run entirely on \code{master} (no network); $p\ge16$ span all three nodes.
All points use the one memory-optimized binary (no swap, mutually
consistent).}\label{tab:d}
\begin{tabular}{@{}rrrrr@{}}
\toprule
$p$ & $T(p)$ no comm (s) & speedup (no comm) & $T(p)$ with comm (s) & speedup (with comm) \\
\midrule
1  & 91.99 & 1.00  & 92.01 & 1.00 \\
2  & 46.58 & 1.97  & 46.62 & 1.97 \\
4  & 24.44 & 3.76  & 25.17 & 3.66 \\
8  & 20.38 & 4.51  & 20.91 & \textbf{4.40} \\
\midrule
16 & 6.16  & 14.93 & 37.53 & 2.45 \\
22 & 4.56  & \textbf{20.15} & 56.75 & 1.62 \\
32 & 4.56  & 20.19 & 30.44 & 3.02 \\
44 & 3.64  & 25.27 & 26.08 & 3.53 \\
\bottomrule
\end{tabular}
\end{table}

\subsection{Chart E --- ring vs.\ 1D tensor path}
\label{sec:charte}
\textbf{Setup.} The same single-head problem at $N=16384$, swept over
$p\in\{2,4,8,16\}$ (\code{scripts/bench\_ring\_vs\_1d.sh}). The left panel is the
measured total time (compute\,+\,comm\,+\,wait); the right panel is the peak number
of $K,V$ rows resident per rank.

\chartfig{results/chart-ring-vs-1d.png}
  {Ring vs.\ 1D: measured time and per-rank $K,V$ memory footprint at $N=16384$.}{fig:e}

\noindent\textbf{Observation --- the ring path wins once the network is real.}
Table~\ref{tab:e} gives the measured split. At $p=2$, where both ranks sit on
\code{master} and communication is essentially free, the ring is slightly slower
($16.5$ vs $14.3$\,s) --- its online-softmax rescaling is pure overhead when there
is no network to hide. But as soon as the job spans nodes the picture inverts: at
$p=16$ the 1D path spends $13.8$\,s in the $K,V$ broadcast and a further $4.2$\,s
waiting at barriers, for a $20.0$\,s total, whereas the ring overlaps its
non-blocking shifts with compute and finishes in $6.3$\,s --- a \textbf{3.2$\times$}
improvement, with communication cut from $13.8$\,s to $3.3$\,s and idle wait from
$4.2$\,s to $0.55$\,s. The right panel shows the structural reason: the 1D path
replicates all $N$ rows of $K,V$ on every rank, while the ring holds only
$\lceil N/p\rceil$ (down to $1024$ rows at $p=16$). The total bytes moved are
comparable; the ring's advantage is \emph{overlap} (latency hiding) and
\emph{memory footprint}, both of which matter precisely on a slow, memory-limited
cluster like this one.

\begin{table}[H]
\centering\small
\caption{Ring vs.\ 1D at $N=16384$ (seconds). Total $=$ compute $+$ comm $+$ wait.
The ring's overlapped, sharded transfers dominate once the run crosses nodes
($p\ge16$).}\label{tab:e}
\begin{tabular}{@{}rcccc@{}}
\toprule
$p$ & 1D comm & 1D total & Ring comm & Ring total \\
\midrule
2  & 0.02 & 14.27 & 0.02 & 16.48 \\
4  & 0.03 & 9.56  & 0.02 & 8.96 \\
8  & 0.03 & 8.25  & 0.04 & 4.96 \\
16 & 13.76 & 19.97 & 3.30 & 6.25 \\
\bottomrule
\end{tabular}
\end{table}

\subsection{Chart F --- OpenMP thread scaling}
\label{sec:chartf}
\textbf{Setup.} One MPI rank on \code{master}, $N=8192$, varying
\code{OMP\_NUM\_THREADS} $\in\{1,2,4,8,16\}$ (\code{scripts/bench\_threads.sh});
pure compute, no communication.

\chartfig{results/chart-threads.png}
  {Intra-node OpenMP runtime and speedup vs.\ thread count ($N=8192$, \code{master}).}{fig:f}

\noindent\textbf{Observation --- a hard memory-bandwidth plateau.} The runtime
drops from $25.8$\,s at one thread to $19.6$\,s at two ($1.32\times$) and then
\emph{does not improve at all} through $4$, $8$ and $16$ threads ($19.4$--$19.5$\,s,
speedup stuck at $\approx\!1.33\times$). The kernel re-reads the entire $K,V$ block
for every query row, so it is \emph{memory-bandwidth bound}: two threads already
saturate \code{master}'s memory channels, and extra threads only contend for the
same bandwidth. This is a genuinely useful negative result --- it explains why the
per-rank (single-thread) MPI decomposition is the right primary lever here, and why
the sharded ring path (whose per-rank $K,V$ block is only $N/p$ rows and stays
closer to cache) is the more scalable route at large $N$ and large $p$. It also
means a hybrid MPI+OpenMP launch buys little on this kernel; the MPI process
dimension, which adds memory channels (and, across nodes, whole separate memory
systems), is what actually scales the compute (cf.\ Chart~D, where crossing from
one node to three lifts the computation-only speedup from $4.5\times$ on $8$ cores
to about $20\times$ on $22$ cores).

% =====================================================================
\section{Discussion: improvements and remaining limitations}
\label{sec:discussion}
% =====================================================================
\paragraph{Improvements implemented.} Four design choices shaped the final system,
each targeting a cost the measurements exposed. The \emph{streaming online-softmax
kernel} processes one query row at a time in $\mathcal{O}(N)$ scratch instead of
forming the $N\times N$ score matrix, lifting the quadratic-memory wall so $N$
reaches $65536$. The \emph{ring (Flash-style) path} replaces the full-$K,V$
broadcast with $\mathcal{O}(N/p)$ sharded blocks streamed by non-blocking shifts
that overlap the compute, cutting cross-node communication $4\times$ and idle wait
to near zero (Chart~E). The \emph{collective hybrid distribution} fans work to the
group roots with \code{Scatter}/\code{Gather} ($\mathcal{O}(\log G)$) rather than a
serial \code{Send}/\code{Recv} loop, and the \emph{rank-0-only input allocation}
relieves non-root ranks of the unused full matrices, removing the memory ceiling on
$p$. An OpenMP layer over the query-row loop adds an intra-node level beneath MPI,
and the built-in profiler separates compute/comm/wait and counts bytes and latency
so every run explains its own time.

\paragraph{Limitations.} Three remain, and each is visible in the data. The kernel
is \emph{memory-bandwidth bound} --- it re-reads $K,V$ for every query row, so
intra-node threading plateaus at $1.33\times$ (Chart~F); cache/query tiling would
raise arithmetic intensity but would forgo the current bit-exact match. The compute
is \emph{portable scalar C} (auto-vectorized at \code{-O2}, not tuned BLAS), so
absolute throughput is modest --- acceptable for a parallelism study. And the
measurements carry honest \emph{noise}: they are sampled at rank~0 (which understates
the slowest rank's comm and, at large $p$, flatters the compute number as rank~0
holds fewer rows), and the Wi-Fi time varies $>2\times$ between repeats
($97$\,s vs $226$\,s for the same $p{=}22$ broadcast), so we draw conclusions from
the qualitative trend rather than any single point. Busy-waiting inside blocking MPI
calls also keeps idle ranks at $100\%$ CPU, which the \code{--profile-wait} timer
makes visible (Chart~C).

% =====================================================================
\section{Conclusion}
% =====================================================================
We built the Transformer multi-head-attention operator from scratch in C\,+\,MPI
and parallelized it three ways --- head (task), tensor (data; 1D row-block, 2D
Cannon, and a ring/Flash-style path), and hybrid (task\,$\times$\,data via
\code{MPI\_Comm\_split}) --- over a single shared streaming kernel, with an OpenMP
layer inside each rank. Every path reproduces the sequential baseline within
tolerance (bit-exact for 1D/head/hybrid, $\sim\!10^{-7}$ for the reordered
Cannon/ring paths), so the parallel program is a faithful solution to the problem,
and the implementation answers each part of the brief --- level of parallelism,
decomposition, mapping, communication, load balancing, and pseudocode --- with a
measured experiment behind every claim.

\paragraph{What the numbers really say.} The honest headline is a \emph{gap between
two kinds of scaling}. The arithmetic divides almost ideally --- computation-only
speedup reaches about $20\times$ on $22$ cores --- but the end-to-end runtime tells
a sterner story: over Wi-Fi the best wall-clock speedup is only $4.4\times$, and it
is reached at $p=8$ on a \emph{single} machine. Bringing in the second and third
nodes makes the computation faster yet the program slower, because the $K,V$
broadcast and gather must cross a link roughly three orders of magnitude slower than
local memory. That same asymmetry turns a statically ``balanced'' hybrid mapping
into a $175.7\%$ load imbalance, which a change of decomposition --- to the
rootless 1D split --- brings down to $28.6\%$.

\paragraph{A realistic appraisal.} Read plainly, \emph{this workload and this
cluster are mismatched}, and the report's worth is in showing precisely why and by
how much. Attention is $\mathcal{O}(N^2)$ compute over $\mathcal{O}(N)$ data, so in
principle it should be compute-bound and scale well --- and it does, on one
multi-core node or whenever the data stays local. The moment that data must travel
over commodity Wi-Fi the constant factors invert: the operator is communication-bound
at every size the $2$--$3$ minute budget permits, and only beyond $N\approx131072$
(a $\sim\!30$-minute run) would computation finally dominate. Distributing attention
across machines is thus worth it precisely when the interconnect is fast (wired
Ethernet or InfiniBand) or the per-node problem is large enough to amortize the
transfer; on a shared wireless LAN the single fastest node is hard to beat. The two
levers that actually move the result are the ones that \emph{reduce and overlap}
data movement (the sharded, non-blocking ring path) and that \emph{place work by
capability} (throughput-proportional, topology-aware distribution on heterogeneous
nodes) --- not the ones that merely divide the arithmetic more finely.

\paragraph{Closing.} The project therefore delivers two things: a correct,
framework-free, three-way-parallel attention, and --- more valuably --- a measured
account of where its time goes and which engineering choices change it. The next
steps follow directly and change none of the parallel structure, only the three
quantities this study isolated. A \emph{wired network} would pull the
with-communication curve toward the computation-only one; \emph{FlashAttention-style
tiling} would lift the memory-bandwidth ceiling that caps intra-node threading; and
\emph{throughput-weighted row assignment} would close the residual heterogeneity
imbalance. The communication wall measured here is not a flaw in the code but a
property of the machine --- and naming it quantitatively is the point of the
exercise.

% =====================================================================
\appendix
\section{Reproducing the results}
% =====================================================================
\begin{enumerate}[leftmargin=1.6em]
\item \textbf{Build:} \code{make} (produces \code{hybrid\_attention}).
\item \textbf{Cluster:} bring up SSH/NFS/OpenMPI and confirm with
      \code{bash scripts/test\_cluster.sh} (wait for \code{cluster OK}); see
      \code{CHECKLIST.md}.
\item \textbf{Correctness:}
      \code{mpirun --hostfile hostfile -np $X$ ./hybrid\_attention --mode hybrid --seq-len 1024}.
\item \textbf{Benchmarks} (set \code{TOTAL\_PROCS} = core count; on our cluster the
      binary is staged to \code{/tmp/mpiattn} on each node and the LAN interface is
      pinned, hence the extra environment variables):
\begin{lstlisting}[numbers=none,mathescape=false]
export BINARY=/tmp/mpiattn/hybrid_attention NET=192.168.124.0/24 HOSTFILE=hostfile
TOTAL_PROCS=22 MODE=hybrid TARGET_S=140 bash scripts/bench_n_sizing.sh    # Chart B
TOTAL_PROCS=22 N=32768 MODE=hybrid       bash scripts/bench_granularity.sh # Chart C
TOTAL_PROCS=44 N=16384 MODE=tensor PLIST="2 4 8 16 22 32 44" \
              bash scripts/bench_speedup.sh                                # Chart D
TOTAL_PROCS=16 N=16384 bash scripts/bench_ring_vs_1d.sh                    # Chart E
MAX_THREADS=16 N=8192  bash scripts/bench_threads.sh                       # Chart F
\end{lstlisting}
\item \textbf{Charts:} \code{python3 scripts/make\_charts.py} (needs
      \code{python3-matplotlib}); writes the five PNGs this report includes.
\end{enumerate}

% =====================================================================
\section{Raw per-process measurements (granularity run)}
\label{app:raw}
% =====================================================================
Tables~\ref{tab:raw-hybrid} and~\ref{tab:raw-tensor} list every process's measured
times for the two granularity runs of Section~\ref{sec:chartc} ($N=32768$, $p=22$,
\code{--profile-wait}), the data behind Chart~C and the $175.7\%\!\to\!28.6\%$
adjustment. ``Work'' $=$ compute $+$ comm is the load-balance quantity; ``wait'' is
idle time at barriers. Note in the hybrid table how the two group roots (ranks~0
and~11) carry far more communication than the members, and how the fast
\code{master} ranks accumulate large \emph{wait} while the communication-bound
\code{slave} ranks accumulate little --- the signature of a communication-driven
imbalance.

\begin{table}[H]
\centering\scriptsize
\caption{Hybrid ($G=2$) per-process times (s). Imbalance on ``work'':
$(147.2-53.4)/53.4=175.7\%$.}\label{tab:raw-hybrid}
\begin{tabular}{@{}rlrrrr@{}}
\toprule
rank & node & compute & comm & wait & work \\ \midrule
0 & \code{master} & 29.0 & 48.0 & 70.7 & 77.0 \\
1 & \code{master} & 35.4 & 18.0 & 83.1 & 53.4 \\
2 & \code{master} & 35.8 & 18.0 & 82.7 & 53.8 \\
3 & \code{master} & 35.7 & 18.0 & 82.8 & 53.7 \\
4 & \code{master} & 35.9 & 18.0 & 82.6 & 53.9 \\
5 & \code{master} & 36.2 & 18.0 & 82.3 & 54.2 \\
6 & \code{master} & 35.6 & 18.0 & 82.8 & 53.7 \\
7 & \code{master} & 35.9 & 18.0 & 82.5 & 54.0 \\
8 & \code{slave1} & 30.8 & 36.8 & 68.9 & 67.6 \\
9 & \code{slave1} & 30.8 & 39.2 & 66.6 & 70.0 \\
10 & \code{slave1} & 29.9 & 40.6 & 66.3 & 70.5 \\
11 & \code{slave1} & 13.8 & 133.4 & 0.3 & \textbf{147.2} \\
12 & \code{slave2} & 12.1 & 73.3 & 50.1 & 85.4 \\
13 & \code{slave2} & 15.7 & 74.8 & 45.1 & 90.5 \\
14 & \code{slave2} & 8.3 & 77.6 & 49.7 & 85.9 \\
15 & \code{slave2} & 11.8 & 79.7 & 44.3 & 91.5 \\
16 & \code{slave2} & 6.5 & 71.1 & 58.2 & 77.7 \\
17 & \code{slave2} & 6.1 & 70.2 & 59.6 & 76.3 \\
18 & \code{slave2} & 13.1 & 53.9 & 69.1 & 67.0 \\
19 & \code{slave2} & 13.9 & 108.6 & 13.7 & 122.5 \\
20 & \code{slave2} & 15.2 & 76.5 & 44.5 & 91.8 \\
21 & \code{slave2} & 12.5 & 58.8 & 65.0 & 71.4 \\
\bottomrule
\end{tabular}
\quad
\begin{tabular}{@{}rlrrrr@{}}
\toprule
rank & node & compute & comm & wait & work \\ \midrule
0 & \code{master} & 4.8 & 69.2 & 5.4 & 74.0 \\
1 & \code{master} & 4.7 & 63.6 & 9.7 & 68.3 \\
2 & \code{master} & 4.7 & 58.8 & 14.5 & \textbf{63.5} \\
3 & \code{master} & 4.8 & 71.3 & 1.8 & 76.2 \\
4 & \code{master} & 4.8 & 68.9 & 4.3 & 73.6 \\
5 & \code{master} & 4.7 & 73.2 & 0.1 & 77.9 \\
6 & \code{master} & 4.7 & 58.8 & 14.4 & 63.5 \\
7 & \code{master} & 5.0 & 71.3 & 1.6 & 76.3 \\
8 & \code{slave1} & 4.4 & 69.7 & 3.8 & 74.1 \\
9 & \code{slave1} & 5.1 & 64.3 & 8.7 & 69.4 \\
10 & \code{slave1} & 5.2 & 63.0 & 10.0 & 68.2 \\
11 & \code{slave1} & 3.5 & 72.0 & 2.7 & 75.5 \\
12 & \code{slave2} & 1.5 & 69.5 & 7.3 & 71.0 \\
13 & \code{slave2} & 1.6 & 74.3 & 2.6 & 75.9 \\
14 & \code{slave2} & 1.7 & 64.1 & 12.8 & 65.7 \\
15 & \code{slave2} & 1.1 & 72.7 & 4.8 & 73.9 \\
16 & \code{slave2} & 2.5 & 68.5 & 7.8 & 70.9 \\
17 & \code{slave2} & 1.8 & 65.5 & 11.6 & 67.3 \\
18 & \code{slave2} & 1.3 & 59.3 & 18.5 & 60.5 \\
19 & \code{slave2} & 1.3 & 72.0 & 5.9 & 73.2 \\
20 & \code{slave2} & 2.2 & 69.6 & 7.4 & 71.8 \\
21 & \code{slave2} & 1.9 & 74.6 & 2.9 & \textbf{77.9} \\
\bottomrule
\end{tabular}
\captionof{table}{1D tensor per-process times (s). Imbalance on ``work'':
$(77.9-60.5)/60.5=28.6\%$.}\label{tab:raw-tensor}
\end{table}

% =====================================================================
\section{Speedup at the doubled size $2N^{\star}=65536$}
\label{app:2n}
% =====================================================================
The brief specifies the speedup study at $2N^{\star}$. We also ran that size; it is
reported here rather than as the headline chart because, as
Section~\ref{sec:chartd} explains, each point exceeds ten minutes of wall time and
the Wi-Fi communication time varies by more than $2\times$ between repeats, so the
$N^{\star}$ sweep gives a cleaner, mutually consistent curve. The qualitative
behaviour at $2N^{\star}$ is identical (Table~\ref{tab:2n}): the computation-only
speedup climbs past $20\times$ while the with-communication speedup is pinned in the
low single digits by the network --- the same communication wall.

\begin{table}[H]
\centering\small
\caption{Speedup at $2N^{\star}=65536$ ($T(1)=401.7$\,s, rank-0 timings).
Single-sampled over Wi-Fi, so the with-communication column carries visible
network noise (e.g.\ the $p=22$ dip from a slow broadcast sample); the trend ---
compute scales, communication caps --- matches Chart~D.}\label{tab:2n}
\begin{tabular}{@{}rrrrr@{}}
\toprule
$p$ & $T(p)$ no comm (s) & speedup (no comm) & $T(p)$ with comm (s) & speedup (with comm) \\
\midrule
1  & 401.65 & 1.00  & 401.69 & 1.00 \\
2  & 210.35 & 1.91  & 210.42 & 1.91 \\
4  & 129.34 & 3.11  & 130.34 & 3.08 \\
8  & 125.45 & 3.20  & 126.65 & 3.17 \\
16 & 31.61  & 12.71 & 107.81 & 3.73 \\
22 & 19.07  & 21.06 & 245.42 & 1.64 \\
32 & 21.42  & 18.75 & 66.69  & 6.02 \\
44 & 23.20  & 17.32 & 74.88  & 5.36 \\
\bottomrule
\end{tabular}
\end{table}

\end{document}
